\documentclass[border=4pt,12pt]{standalone}
% House style for redrawn exam diagrams (diagram-v4: 0.4pt lines, filled -latex arrows, fine dashes; stipple and hatching drawn as real dots/lines (no PDF/SVG patterns, so they render in every browser and at any scale); transparent background). Every figure \input's this.
\usepackage{amsmath,amssymb}
\usepackage{txfonts}
\usepackage{pdfrender,xcolor}
% \halo{text}: label on shading. The shading stops just around each letter's outline (a thin knock-out that follows the
% glyphs, like the papers' typesetting), never a box. Owner: "I don't want a box around the text".
\newcommand{\halo}[1]{\leavevmode\rlap{\pdfrender{TextRenderingMode=2,LineWidth=2.2pt,LineJoinStyle=1,StrokeColor=white,FillColor=white}#1}#1}
\usepackage{tikz,pgfplots}
\pgfplotsset{compat=1.18}
\usetikzlibrary{arrows.meta,calc,angles,quotes,patterns,patterns.meta,intersections,positioning,decorations.markings,shapes.geometric}
% Stipple and hatching are drawn as actual dots / lines clipped to the region (path picture), NOT as pattern fills:
% SVG <pattern> fills vanish or blur when browsers scale the image, so the owner could not see the shading.
% Each dot is a 0.72pt square (keeps SVGs small; reads as a dot at print size). Dots sit on one global staggered grid, so neighbouring regions line up.
\makeatletter
\newdimen\ex@x \newdimen\ex@y \newdimen\ex@xa \newdimen\ex@ya \newdimen\ex@xb \newdimen\ex@yb \newcount\ex@row
\def\exam@stipplecolor{black!75}\def\exam@dotr{0.42pt}\def\exam@dotgap{2.5}
\def\ex@rowdots{\ifdim\ex@x<\ex@xb \pgfpathrectangle{\pgfqpoint{\dimexpr\ex@x-0.36pt\relax}{\dimexpr\ex@y-0.36pt\relax}}{\pgfqpoint{0.72pt}{0.72pt}}\advance\ex@x by \exam@dotgap pt\relax\expandafter\ex@rowdots\fi}
\def\exam@stipple{%
  \pgftransformreset% dots are sized in real points even inside scaled pictures
  \pgfpointanchor{path picture bounding box}{south west}\ex@xa=\pgf@x \ex@ya=\pgf@y
  \pgfpointanchor{path picture bounding box}{north east}\ex@xb=\pgf@x \ex@yb=\pgf@y
  \pgfsetfillcolor{\exam@stipplecolor}\pgfmathsetmacro\ex@h{\exam@dotgap/2}%
  \pgfmathtruncatemacro\ex@r{floor(\ex@ya/\ex@h)-1}\ex@row=\ex@r\relax
  \ex@y=\ex@r\dimexpr\ex@h pt\relax
  \loop\ifdim\ex@y<\ex@yb
    \pgfmathtruncatemacro\ex@c{floor(\ex@xa/\exam@dotgap)-1}\ex@x=\ex@c\dimexpr\exam@dotgap pt\relax
    \ifodd\ex@row \advance\ex@x by \ex@h pt\fi
    \ex@rowdots
    \advance\ex@y by \ex@h pt \advance\ex@row by 1
    \ifnum\numexpr\ex@row-(\ex@row/4)*4\relax=0 \pgfusepath{fill}\fi% flush every 4 rows: bounded TeX memory, few clip copies in the SVG
  \repeat\pgfusepath{fill}}
\newdimen\ex@d
\def\exam@hatch#1{% #1 = 1 for /, -1 for \ ; lines 45 degrees, 3pt apart
  \pgftransformreset
  \pgfpointanchor{path picture bounding box}{south west}\ex@xa=\pgf@x \ex@ya=\pgf@y
  \pgfpointanchor{path picture bounding box}{north east}\ex@xb=\pgf@x \ex@yb=\pgf@y
  \ex@d=\dimexpr\ex@yb-\ex@ya\relax \ex@x=\dimexpr\ex@xa-\ex@d\relax
  \loop\ifdim\ex@x<\ex@xb
    \ifnum#1>0 \pgfpathmoveto{\pgfqpoint{\ex@x}{\ex@ya}}\pgfpathlineto{\pgfqpoint{\dimexpr\ex@x+\ex@d\relax}{\ex@yb}}%
    \else \pgfpathmoveto{\pgfqpoint{\ex@x}{\ex@yb}}\pgfpathlineto{\pgfqpoint{\dimexpr\ex@x+\ex@d\relax}{\ex@ya}}\fi
    \advance\ex@x by 4.24pt
  \repeat
  \pgfsetlinewidth{0.4pt}\pgfsetstrokecolor{black}\pgfusepath{stroke}}
\def\examstipple{\exam@stipple}\def\examhatch{\exam@hatch}\def\examsetstipplecolor#1{\def\exam@stipplecolor{#1}}\def\examsetdotgap#1{\def\exam@dotgap{#1}}
\makeatother
% (legacy pattern kept for old files)
% staggered stipple like the NESA papers' dotted shading
\pgfdeclarepatternformonly{exam stipple}{\pgfpointorigin}{\pgfqpoint{2.2pt}{2.2pt}}{\pgfqpoint{2.2pt}{2.2pt}}{%
  \pgfpathcircle{\pgfqpoint{0.55pt}{0.55pt}}{0.42pt}\pgfpathcircle{\pgfqpoint{1.65pt}{1.65pt}}{0.42pt}\pgfusepath{fill}}
\tikzset{
  every picture/.style={line width=0.6pt, >={Latex[length=2mm,width=1.3mm]}},
  curve/.style={line width=0.6pt},
  open arrow/.style={-{Latex[length=2mm,width=1.3mm]}},  % owner: filled -latex heads everywhere, no open chevrons
  aux/.style={dash pattern=on 2pt off 1.5pt, line width=0.6pt},
  shade/.code={\def\exA{#1}\def\exB{none}\ifx\exA\exB\else\pgfkeysalso{fill=black!18}\fi},
  shade/.default=on,
  shade dots/.style={fill=none, path picture={\examstipple}},
  stipple color/.code={\examsetstipplecolor{#1}},
  stipple gap/.code={\examsetdotgap{#1}},
  shade dense dots/.style={fill=none, stipple gap=1.6, path picture={\examstipple}},
  shade hatch/.style={fill=none, path picture={\examhatch{1}}},
  shade hatch reverse/.style={fill=none, path picture={\examhatch{-1}}},
  shade solid/.style={fill=black},
  pt/.style={circle, fill, inner sep=1.1pt},
  open pt/.style={circle, draw, fill=white, inner sep=1.1pt},
}
\pgfplotsset{
  exam axes/.style={axis lines=middle, axis line style={->}, xlabel={$x$}, ylabel={$y$},
    xlabel style={anchor=west}, ylabel style={anchor=south}, tick label style={font=\small},
    every axis plot/.append style={line width=0.6pt, black}, samples=200, clip=false,
    enlargelimits=false, axis on top},
}

\begin{document}
% 159_qVI4_s1: NOT TO SCALE. y = f(x) as in 159_qVI4_1 (same traced curve, placed at this scan's position, origin
% (696,486)); y = g(x) = 3f(x-2) built from that curve: x mapped so f's points -3, -1, 2, 5, 7 land where the teacher
% put g's -1, 1, 4, 7, 9 (smooth monotone map), y stretched by one factor about the x-axis (3.75 px per f-px), so
% g's turning points and intercepts are the exact images of f's (rule 16) and the curve is smooth (rule 10).
\begin{tikzpicture}[x=0.0127cm,y=-0.0127cm]
  \draw[<->] (18,486) -- (1372,486);
  \node[anchor=base west, inner sep=0pt] at (1380,494) {$x$};
  \draw[<->] (696,922) -- (696,50);
  \node[anchor=base, inner sep=0pt] at (696,32) {$y$};
  \node[align=left, anchor=north west, inner sep=0pt] at (124,55) {NOT TO\\SCALE};
  \draw[<->, line width=0.9pt] plot[smooth] coordinates {(292.5,917) (300,885.7) (307.6,855) (315.5,824.3) (323.7,794) (332.3,764.1) (341.3,734.4) (350.9,704.9) (361.1,676) (371.8,647.4) (383.3,619.3) (395.4,591.4) (408.3,563.8) (421.7,537.4) (435.6,512.8) (450.4,490.1) (466.1,469.8) (482.9,452.2) (500.3,437.7) (519,426) (538.4,417.7) (558.7,412.3) (580,410) (601,410.9) (623.5,415) (646.5,422.2) (669.9,432.2) (693.7,444.8) (717.7,459.7) (742.5,476.7) (768.3,495) (793.5,512.4) (817.6,527.7) (840.9,540.5) (863.9,550.5) (886.1,557.5) (908.2,561.5) (929.8,562.4) (950.4,560) (970.4,554.4) (990.5,545.5) (1010,533.6) (1028,519) (1045,501.7) (1061,481.9) (1076,459.8) (1089,435.7) (1102,409.8) (1114,382.6) (1126,354.2) (1137,325.3) (1148,296.2) (1159,266.8) (1169,237) (1178,207.5) (1187,177.4) (1196,147) (1204,116.6) (1213,86.08) (1220,55)};
  \draw[<->, line width=0.9pt] plot[smooth] coordinates {(596.3,830) (599.2,790.3) (601.9,750.6) (604.5,710.4) (607,670.7) (609.7,631) (612.5,591.3) (615.8,551.6) (619.5,512.6) (623.8,473.6) (628.7,435.2) (634.5,396.9) (641.3,358.7) (649.1,320.7) (657.7,284.4) (668.1,252.4) (681.4,226.6) (697.8,206.6) (712.6,197.3) (725.4,196.1) (740.9,202.4) (757.4,217.9) (776.5,245) (794.4,276.3) (811.2,310.1) (827.7,345.5) (843.1,381) (858.2,417.4) (872.8,454) (886.9,490.8) (901.3,528.5) (915.4,565.6) (929.7,602) (944.6,637.6) (960.4,672.3) (977.3,705.6) (994.6,735.3) (1013,757) (1032,770.7) (1051,776.4) (1072,774.7) (1094,765.4) (1115,748.5) (1136,724.7) (1156,696.7) (1172,666.5) (1186,634.1) (1198,599.7) (1207,564.5) (1215,527.6) (1221,490.2) (1226,452) (1230,412.7) (1234,373) (1238,333.5) (1242,293.5) (1246,254.1) (1249,214.6) (1252,174.9) (1255,135)};
  % f
  \foreach \p in {(455,486),(755,486),(1055,486),(582,408.5),(930,564)} \fill \p circle (4.5);
  \node[anchor=north, inner sep=0pt] at (464,500) {$-3$};
  \node[anchor=north, inner sep=0pt] at (762,500) {$2$};
  \node[anchor=north, inner sep=0pt] at (1062,500) {$7$};
  \node[anchor=south, inner sep=0pt] at (568,397) {$(-1,3)$};
  \node[anchor=north west, inner sep=0pt] at (942,580) {$(5,-4)$};
  % g = 3f(x-2)
  \foreach \p in {(623,486),(885,486),(1221,486),(720,195.4),(1057,778.5)} \fill \p circle (4.5);
  \node[anchor=north east, inner sep=0pt] at (604,500) {$-1$};
  \node[anchor=north east, inner sep=0pt] at (879,500) {$4$};
  \node[anchor=north west, inner sep=0pt] at (1228,500) {$9$};
  \node[anchor=south west, inner sep=0pt] at (733,180) {$(1,9)$};
  \node[anchor=north, inner sep=0pt] at (1057,800) {$(7,-12)$};
\end{tikzpicture}
\end{document}
